\documentclass[11pt]{article}
\usepackage[T1]{fontenc}
\usepackage[utf8]{inputenc}
\usepackage{lmodern}
\usepackage{amsmath,amssymb}
\usepackage{longtable,booktabs,array}
% Compatibility with Pandoc's captionless tables on older longtable versions.
\ifcsname c@none\endcsname\else\newcounter{none}\fi
\usepackage{graphicx}
\usepackage[margin=25mm]{geometry}
\usepackage{hyperref}
\hypersetup{colorlinks=true,linkcolor=blue,urlcolor=blue}
\providecommand{\tightlist}{\setlength{\itemsep}{0pt}\setlength{\parskip}{0pt}}
\providecommand{\pandocbounded}[1]{#1}
\setlength{\emergencystretch}{3em}
\setcounter{secnumdepth}{-1}
\begin{document}
\section{The theorem on formal
functions}\label{the-theorem-on-formal-functions}

\emph{Written by GPT-6.1 Sol (OpenAI), in Codex, at Ultra effort,
October 2026. Self-checked by the writing AI, GPT-6.1 Sol, at Ultra
effort. Independent portions are public domain (CC0). Appendix Z adapts
the Stacks project authors' free conductor and strong-transcendence
proofs and retains GFDL-1.2-or-later. See the
\href{licensing.html}{course notice} for authorship and component
terms.}

A fiber remembers a family at one parameter. Its successive
infinitesimal neighborhoods remember powers of that parameter's maximal
ideal. Formal functions identifies the inverse limit of their cohomology
with the completion of a finite cohomology module. The individual
comparison maps can fail to be isomorphisms; what makes the theorem work
is a uniform bound on how long their errors persist.

We use \href{proper-morphisms-and-coherent-direct-images.md}{proper
coherent direct images},
\href{affine-cohomology-and-serres-criterion.md}{affine cohomology and
Serre's criterion}, and
\href{base-change-and-the-grothendieck-complex.md}{flat base change}.
\href{../AG-CA/AG-CA-19.html\#3-noetherian-exactness-tensoring-and-flatness}{Completion,
Theorems 3.1--3.3} supplies exact completion of finite modules over
Noetherian rings and faithful flatness when the completion ideal lies in
the Jacobson radical, including maximal-ideal completion of a Noetherian
local ring. The dimension consequences in Section 5 use
\href{../ag-etale-cohomology/cohomological-dimension-and-the-kunneth-formula.html\#7-the-bound-for-all-finite-type-schemes}{Cohomological
dimension and the Künneth formula, Lemma 7.1}, the full Noetherian
topological vanishing proof for arbitrary abelian sheaves. No flatness
assumption on the coherent sheaf is imposed in this lesson.

\subsection{0. A normality input in the algebraic Zariski Main
proof}\label{a-normality-input-in-the-algebraic-zariski-main-proof}

The quasi-finite input used in Section 5 has an algebraic Zariski Main
proof whose supporting coefficient argument requires polynomial
normality over an arbitrary normal domain. A proof for polynomial rings
over fields alone would not suffice. We supply the general statement
here. The valuation existence input is the written
\href{../AG-GS/prerequisites/AG-MO/valuation-rings-and-separatedness.html}{Valuation
rings and separatedness, Theorem 2.1}; it gives a valuation ring in a
specified field dominating a specified local subring, without a
Noetherian hypothesis or a discrete-value assumption.

\textbf{Lemma 0.1 (polynomial normality).} If \(R\) is an integrally
closed domain, then \(R[x_1,\ldots,x_n]\) is an integrally closed domain
for every finite \(n\).

\textbf{Proof.} Put \(F=\operatorname{Frac}R\). First we prove that
\(R\) is the intersection of all valuation subrings of \(F\) containing
\(R\). If \(z\in F\setminus R\), then \(z\) is not integral over \(R\).
In the ring \(R[z^{-1}]\), the element \(z^{-1}\) is not a unit: an
equation \[
1=z^{-1}(a_0+a_1z^{-1}+\cdots+a_mz^{-m})
\] would yield a monic polynomial equation for \(z\) over \(R\), after
multiplication by \(z^{m+1}\). Choose a maximal ideal containing
\(z^{-1}\) and localize at it. The earlier valuation theorem gives a
valuation ring \(V\subset F\) dominating this local ring. Thus
\(R\subset V\), while \(z^{-1}\) is in its maximal ideal and
\(z\notin V\). This separates every \(z\notin R\) from the intersection
and proves the assertion.

For such a valuation ring, let \(v:F^*\to\Gamma\) be its ordered value
group, with \(V=\{0\}\cup\{a:v(a)\geq0\}\). This group can be
constructed as \(F^*/V^*\), ordered by divisibility in \(V\). For a
nonzero polynomial \(h=\sum a_ix^i\in F[x]\), define \[
w(h)=\min_i v(a_i),
\] ignoring zero coefficients. We have \(w(hg)=w(h)+w(g)\): divide each
polynomial by a coefficient of least value. The resulting polynomials
have coefficients in \(V\) and nonzero reductions in
\((V/\mathfrak m_V)[x]\); their product has nonzero reduction because
this polynomial ring over a field is a domain. Also
\(w(h+g)\geq\min(w(h),w(g))\), with equality when the two values differ.
Consequently a sum with a unique term of strictly least value cannot
vanish.

Let \(u\in F(x)\) be integral over \(R[x]\). It is integral over
\(F[x]\), so \(u\in F[x]\). To justify this last assertion directly,
write \(u=a/b\) with coprime polynomials in the Euclidean domain
\(F[x]\). A monic integral equation implies that \(b\) divides \(a^m\).
Bézout's identity for \(a,b\) forces \(b\) to be a unit. Thus \(u\) is a
polynomial.

For every valuation ring \(V\) above, write its integral equation as \[
u^m+c_1(x)u^{m-1}+\cdots+c_m(x)=0,
\qquad c_i(x)\in R[x]\subset V[x].
\] If \(w(u)<0\), the term \(u^m\) has value \(m w(u)\), strictly less
than the value of every other nonzero term, since \(w(c_i)\geq0\). This
contradicts the least-value property. Hence every coefficient of \(u\)
belongs to \(V\). The intersection assertion puts all coefficients in
\(R\), proving that \(R[x]\) is integrally closed. Induction on the
number of variables proves the lemma. \(\square\)

This proof supplies the polynomial-normality step of the supporting
algebraic argument. It does not certify all its other inputs; their
exact earlier proof locations and remaining audit boundaries are
recorded in the prerequisite record.

\subsection{1. The completion map}\label{the-completion-map}

Let \(A\) be Noetherian, \(I\subset A\) an ideal,
\(f:X\to\operatorname{Spec}A\) proper, and \(F\) coherent. For
\(n\geq1\), put \[
X_n=X\times_A\operatorname{Spec}(A/I^n),\qquad F_n=F|_{X_n}.
\] The closed immersion \(i_n:X_n\hookrightarrow X\) has
\(i_{n*}F_n=F/I^nF\). Since closed immersion pushforward is exact and
has no positive direct images, we may compute \[
H^q(X_n,F_n)=H^q(X,F/I^nF).
\] Set \(M^q=H^q(X,F)\), a finite \(A\)-module by proper finiteness. The
quotient map on sheaves induces \[
M^q/I^nM^q\longrightarrow H^q(X_n,F_n),
\] because \(I^n\) annihilates the target. These maps are compatible
with reduction from \(n+1\) to \(n\). We seek an isomorphism from \[
\widehat{M^q}=\varprojlim_n M^q/I^nM^q
\] to the inverse limit of the right-hand terms. Completion here is with
respect to \(I\); neither \(A\) nor \(M^q\) is assumed complete in
advance.

The theorem is not a claim that every finite-level map is an
isomorphism. For example, the extension bundle on \(\mathbf P^1_{k[t]}\)
with class \(t\) has \(H^0(X,E)=0\), but has a nonzero section on the
fiber \(t=0\). We will return to its infinitesimal neighborhoods in
Exercise 7.6.

\subsection{2. All powers at once}\label{all-powers-at-once}

The \textbf{Rees algebra} is \[
R=\bigoplus_{n\geq0}I^nT^n\subset A[T],\qquad R_0=A.
\] If \(I=(a_1,\ldots,a_r)\), then \(R\) is generated by the degree-one
elements \(a_iT\), hence is a quotient of a polynomial algebra in
finitely many variables over \(A\). It is Noetherian. The symbol \(T\)
records degree; it need not itself belong to \(R\).

\textbf{Lemma 2.1 (graded cohomological finiteness).} For each
\(q\geq0\), the graded module \[
Q^q=\bigoplus_{n\geq0}H^q(X,I^nF)T^n
\] is finite over \(R\).

\textbf{Proof.} Form \(X_R=X\times_A\operatorname{Spec}R\), with affine
projection \(\pi:X_R\to X\). On \(X\), the graded sheaf
\(\bigoplus I^nF\) is a module over \(\mathcal O_X\otimes_A R\): a
degree-\(d\) element acts by multiplication into the degree shifted by
\(d\). It is a quotient of \(F\otimes_A R\), by the degreewise maps
\(F\otimes_A I^n\to I^nF\). The affine module/sheaf correspondence gives
a coherent sheaf \(G\) on \(X_R\) with \[
\pi_*G=\bigoplus_{n\geq0}I^nF.
\] Coherence follows because \(X_R\) is Noetherian and this sheaf is
locally a finite module, generated by a finite set of degree-zero
generators of \(F\).

The map \(X_R\to\operatorname{Spec}R\) is proper by base change. Proper
cohomology finiteness makes \(H^q(X_R,G)\) finite over \(R\). Affine
pushforward and Leray identify it with \(H^q(X,\pi_*G)\). Finally
cohomology of quasi-coherent sheaves on the quasi-compact,
quasi-separated scheme \(X\) commutes with filtered colimits, by the
third lesson. A direct sum is the filtered colimit of finite partial
sums, so \[
H^q(X,\pi_*G)=\bigoplus_{n\geq0}H^q(X,I^nF).
\] The identifications preserve multiplication and grading. This is the
asserted finite graded module. \(\square\)

We next extract two consequences. Write \[
J_n^q=\operatorname{im}\bigl(H^q(X,I^nF)\to M^q\bigr),
\qquad K_n^q=\ker\bigl(H^q(X,I^nF)\to M^q\bigr).
\] Both collections respect graded multiplication. Indeed inclusion of
ideal powers commutes with multiplying a section or cohomology class by
an element of \(I^d\). Thus \(\bigoplus J_n^qT^n\) is a finite graded
image module of \(Q^q\), and \(\bigoplus K_n^qT^n\) is a finite graded
submodule, using Noetherianness of \(R\).

\textbf{Lemma 2.2 (bounds on the image and kernel).} There are integers
\(c_J,c_K\geq0\) such that \[
I^nM^q\subset J_n^q\subset I^{n-c_J}M^q\quad(n\geq c_J),
\] and the transition map \(K_m^q\to K_n^q\) is zero whenever
\(m\geq n+c_K\).

\textbf{Proof.} Multiplication by an element of \(I^n\) factors as
\(F\to I^nF\to F\), proving the first inclusion. Choose finitely many
homogeneous generators of the image module, all in degrees at most
\(c_J\). The degree-\(n\) part is a sum of \(I^{n-d}J_d^q\) over
generator degrees \(d\leq c_J\). Each is contained in \(I^{n-c_J}M^q\),
proving the second inclusion. This is the cohomological analogue of an
Artin--Rees bound.

Choose homogeneous kernel generators \(k_\ell\in K_{d_\ell}^q\) with
\(d_\ell\leq c_K\). An element of \(K_m^q\) is a sum of products
\(a_\ell k_\ell\), with \(a_\ell\in I^{m-d_\ell}\). If \(m\geq n+c_K\),
each such coefficient is a sum of products \(ba'\), where \(b\in I^n\)
and \(a'\in I^{m-d_\ell-n}\). Under transition to \(H^q(X,I^nF)\), its
product with \(k_\ell\) is computed by first including \(I^{d_\ell}F\)
into \(F\), then multiplying by \(a'\), then by \(b\) into \(I^nF\). The
initial inclusion sends \(k_\ell\) to zero, by definition of the kernel.
Hence every product maps to zero. This proves the uniform transition
bound. \(\square\)

A system whose sufficiently distant transition maps to each fixed term
are zero is called \textbf{pro-zero}. Its individual terms may be
nonzero; its inverse limit is zero because each coordinate of a
compatible family is the image of a distant, necessarily zero,
transition. Lemma 2.2 supplies this stronger, bounded form of
disappearance for the kernel system.

\subsection{3. Stabilization of the thickened
cohomology}\label{stabilization-of-the-thickened-cohomology}

The short exact sequence \(0\to I^nF\to F\to F/I^nF\to0\) gives the
exact sequence \[
0\to M^q/J_n^q\longrightarrow H^q(X_n,F_n)\longrightarrow K_n^{q+1}\to0.
\tag{1}
\] It is a sequence of inverse systems. The left transitions are
surjective because \(J_m^q\subset J_n^q\) for \(m\geq n\). The right
system is pro-zero by Lemma 2.2 in degree \(q+1\).

\textbf{Proposition 3.1 (the Mittag--Leffler bound).} For fixed \(q\),
choose \(c\) as the kernel bound in degree \(q+1\). Then for
\(m\geq n+c\), \[
\operatorname{im}\bigl(H^q(X_m,F_m)\to H^q(X_n,F_n)\bigr)
=\operatorname{im}\bigl(M^q\to H^q(X_n,F_n)\bigr).
\] Consequently the thickened-cohomology system is Mittag--Leffler.

\textbf{Proof.} In diagram (1) for indices \(m,n\), the right transition
is zero. Therefore the middle transition lands in the left submodule
\(M^q/J_n^q\). It contains all of that submodule: the left transition
\(M^q/J_m^q\to M^q/J_n^q\) is onto, and these left modules inject into
the respective middle terms. Thus its image is exactly the claimed
submodule. The image is independent of all sufficiently large \(m\),
which is the Mittag--Leffler condition. \(\square\)

We can now take the limit without an unexplained exchange of limits and
cohomology. If \((u_n)\) is a compatible family of middle classes in
(1), its images in \(K_n^{q+1}\) form a compatible family in a pro-zero
system, so are all zero. Thus each \(u_n\) lies in the injective left
term, with a unique preimage. These preimages are automatically
compatible. We obtain the canonical identification \[
\varprojlim_n H^q(X_n,F_n)=\varprojlim_n M^q/J_n^q.
\tag{2}
\] This also proves why a transient class on a small thickening need not
define a class in the inverse limit. Compatibility with every higher
thickening is a real constraint.

\subsection{4. Global and stalk forms of formal
functions}\label{global-and-stalk-forms-of-formal-functions}

\textbf{Theorem 4.1 (formal functions).} Under the hypotheses of Section
1, the canonical maps give an isomorphism \[
\widehat{H^q(X,F)}\xrightarrow{\sim}\varprojlim_n H^q(X_n,F_n).
\] It is an isomorphism of \(\widehat A\)-modules and a homeomorphism
for their inverse-limit topologies.

\textbf{Proof.} Combine (2) with the inclusions from Lemma 2.2: \[
I^nM^q\subset J_n^q,\qquad J_{n+c_J}^q\subset I^nM^q.
\] They say that the two descending filtrations \((I^nM^q)\) and
\((J_n^q)\) are cofinal. The identification (2) also preserves the limit
topologies: each left term is a discrete submodule of the corresponding
middle term, and its unique coordinate preimage is continuous on that
submodule. The first inclusions give the canonical continuous map
between the inverse limits of the corresponding quotients. The second
give its inverse by taking the cofinal shifted indices \(n+c_J\). The
two composites are the usual transitions and hence induce the identity
on limits. Both maps are continuous for the quotient limit topologies,
proving the homeomorphism. The maps arise from \(A\)-linear maps at all
levels, where \(I^n\) annihilates the cohomology of \(F_n\), and so
preserve the induced \(\widehat A\)-action. \(\square\)

\textbf{Corollary 4.2 (the stalk form).} Let \(f:X\to S\) be proper,
\(S\) locally Noetherian, \(F\) coherent, and \(s\in S\). Put
\(A_s=\mathcal O_{S,s}\), with maximal ideal \(\mathfrak m_s\), and \[
X_n=X\times_S\operatorname{Spec}(A_s/\mathfrak m_s^n),\qquad F_n=F|_{X_n}.
\] Then \[
\widehat{(R^qf_*F)_s}\cong\varprojlim_n H^q(X_n,F_n)
\] as \(\widehat{A_s}\)-modules.

\textbf{Proof.} The canonical map \(\operatorname{Spec}A_s\to S\) is
flat: on an affine neighborhood it is localization. Flat base change and
the affine description of quasi-coherent direct images identify the
stalk with
\(H^q(X\times_S\operatorname{Spec}A_s,F|_{X\times_S\operatorname{Spec}A_s})\).
The local ring is Noetherian, the changed morphism proper, and the
changed sheaf coherent. Apply Theorem 4.1 with ideal \(\mathfrak m_s\).
Its thickenings are precisely the displayed \(X_n\), since their
parameter maps factor through \(\operatorname{Spec}A_s\). \(\square\)

\subsection{5. Dimension and finiteness
consequences}\label{dimension-and-finiteness-consequences}

\textbf{Theorem 5.1 (vanishing above a fiber's dimension).} For the
morphism of Corollary 4.2, \[
(R^qf_*F)_s=0\quad\text{if }q>\dim X_s.
\]

\textbf{Proof.} Each \(X_n\) has the same underlying topological space
as \(X_s\): the extra ideal is nilpotent and therefore belongs to every
prime. That space is Noetherian, of finite dimension \(d\), because
\(X_s\) is of finite type over a field. The Noetherian topological
vanishing theorem,
\href{../ag-etale-cohomology/cohomological-dimension-and-the-kunneth-formula.html\#7-the-bound-for-all-finite-type-schemes}{Cohomological
dimension and the Künneth formula, Lemma 7.1}, gives \(H^q(X_n,F_n)=0\)
for \(q>d\). Formal functions makes the completed stalk zero. The
uncompleted stalk is finite by proper coherence. For a finite module
over a Noetherian local ring, completion preserves the residue-field
quotient; zero completion implies \(M/\mathfrak mM=0\), and Nakayama
implies \(M=0\). This proves the assertion. An empty fiber gives empty
thickenings and the same conclusion in every degree. \(\square\)

The vanishing is a statement about the whole stalk, stronger than a
statement about its tensor with the residue field. In particular, if all
fiber dimensions are at most \(d\), then \(R^qf_*F=0\) for \(q>d\),
without a flatness assumption.

\textbf{Theorem 5.2 (proper with finite fibers is finite).} A proper
morphism over a locally Noetherian base whose fibers have finitely many
points is finite.

\textbf{Proof.} Its fibers are finite type over fields. Such schemes
with finitely many points are zero-dimensional by Noether normalization.
Indeed an affine chart admits a finite integral surjection onto affine
space of its dimension; a positive-dimensional affine space over any
field has infinitely many points, so a finite-point chart cannot have
positive dimension. Their zero-dimensional coordinate rings are
finite-dimensional by Noether normalization. Thus Theorem 5.1 makes
\(R^1f_*G=0\) for every coherent \(G\).

Work over an affine Noetherian neighborhood \(S=\operatorname{Spec}A\).
Its inverse image \(X\) is Noetherian, separated and quasi-compact. Any
finite type ideal \(J\subset\mathcal O_X\) is coherent, and Leray over
the affine base gives \[
H^1(X,J)=\Gamma(S,R^1f_*J)=0.
\] Here affine vanishing eliminates the other Leray terms, since the
direct images are quasi-coherent. The finite-ideal version of Serre's
affine criterion from the third lesson now implies that \(X\) is affine.
Its function ring \(H^0(X,\mathcal O_X)\) is a finite \(A\)-module by
proper cohomology finiteness. Hence \(X\to\operatorname{Spec}A\) is
finite. This is local on the base and proves the theorem. \(\square\)

A quasi-finite proper morphism over a locally Noetherian base is
therefore finite. The existing earlier \emph{Zariski's Main Theorem}
lesson in \emph{Morphisms of schemes} gives another route; the present
route uses formal functions and the cohomological affine criterion.

\textbf{Corollary 5.3 (a finite fiber has a finite neighborhood).} If
\(f:X\to S\) is proper, \(S\) is locally Noetherian, and \(X_s\) is a
finite set, then \(f\) is finite over an open neighborhood of \(s\).

\textbf{Proof.} The quasi-finite locus is open by the existing earlier
\href{../AG-MO/AG-MO-12.html\#3-from-local-equalities-to-finite-affine-completions}{\emph{Zariski's
Main Theorem}, Theorem 3.1}. It contains the entire finite fiber: the
pointwise finite-fiber test is proved in \emph{Quasi-finite morphisms
and Chevalley}, Theorem 1.1. Its closed complement has closed image in
\(S\), by properness, and that image misses \(s\). Remove the image. The
restricted map has only finite fibers, and Theorem 5.2 makes it finite.
\(\square\)

\textbf{Lemma 5.4 (cartesian nilpotent thickenings).} Let
\(Y_0\hookrightarrow Y\) be defined by an ideal \(J\) with \(J^N=0\),
let \(f:X\to Y\), and put \(X_0=X\times_Y Y_0\). If \(f_0:X_0\to Y_0\)
is of finite type, a closed immersion, or proper, respectively, then so
is \(f\).

\textbf{Proof.} First a nilpotent thickening \(T\) of an affine scheme
\(T_0\) is affine. Their topological spaces coincide, so \(T\) is
quasi-compact and quasi-separated. For every quasi-coherent \(M\) on
\(T\), its finite filtration by powers of the nilpotent ideal has layers
which are quasi-coherent on \(T_0\). Affine acyclicity and the
cohomology sequences give \(H^1(T,M)=0\). The Serre affine criterion
proved in the affine-cohomology lesson now gives affineness.

Work over an affine of \(Y\), and cover \(X_0\) by finitely many affines
if \(f_0\) is of finite type. The corresponding open thickenings in
\(X\) are affine by the preceding argument. On one of them, write the
ring map \(B\to C\). If \(C/JC\) is generated over \(B/J\) by finitely
many elements, lift them and let \(D\subset C\) be the \(B\)-subalgebra
they generate. Then \(C=D+JC\); iteration gives \(C=D+J^NC=D\). This
proves finite type. If \(f_0\) is a closed immersion, \(X_0\) over this
affine is affine, hence \(X\) is affine. Surjectivity modulo \(J\) gives
\(C=\operatorname{im}(B)+JC\); the same iteration makes \(B\to C\)
surjective, proving closed immersion. These affine conclusions are local
on the target.

If \(f_0\) is proper, finite type was just proved. Its diagonal is a
closed immersion; the diagonal of \(f_0\) is the cartesian
nilpotent-base restriction of the diagonal of \(f\), so the
closed-immersion result makes \(f\) separated. Finally, after every base
change, both the source and target have the same underlying spaces as
their nilpotent-base restrictions. Images of closed subsets are
therefore closed, because the base change of \(f_0\) is closed. This
proves universal closedness and hence properness. \(\square\)

\subsection{6. The blow-up of a plane
point}\label{the-blow-up-of-a-plane-point}

Let \(b:X\to\mathbf A^2_k\) be the blow-up of the origin and
\(E\cong\mathbf P^1_k\) its exceptional curve. Its two charts are \[
U=\operatorname{Spec}k[x,v],\quad y=xv,
\qquad V=\operatorname{Spec}k[u,y],\quad x=uy.
\] They identify the blow-up with the closed incidence scheme in
\(\mathbf A^2\times\mathbf P^1\), so \(b\) is projective and proper. Off
the origin it is an isomorphism. The pulled-back origin ideal
\(\mathfrak m=(x,y)\) is generated by \(x\) on \(U\) and by \(y\) on
\(V\), and is the invertible ideal \(J=\mathcal O_X(-E)\). On the
overlap \(y=vx\), so its conormal restriction is \(\mathcal O_E(1)\).
Consequently \[
J^j/J^{j+1}\cong\mathcal O_E(j)\quad(j\geq0).
\] The positive sign is essential: the normal bundle of \(E\) is
\(\mathcal O_E(-1)\), whereas the successive ideal layers use its dual.

The thickened fiber \(X_n\) is defined by \(J^n\). For \(n\geq2\), \[
0\to\mathcal O_E(n-1)\to\mathcal O_{X_n}\to\mathcal O_{X_{n-1}}\to0.
\] Projective-line cohomology gives \(H^1(E,\mathcal O_E(j))=0\) for
every \(j\geq0\). The long exact sequences, starting with \(X_1=E\),
give \(H^1(X_n,\mathcal O_{X_n})=0\) for all \(n\). By formal functions
and Nakayama, \((R^1b_*\mathcal O_X)_0=0\). Away from zero the morphism
is an isomorphism, so \[
R^1b_*\mathcal O_X=0.
\] All still higher direct images vanish by the fiber-dimension bound,
because the largest fiber dimension is one.

We can also see the completed degree-zero comparison explicitly. The
canonical polynomial map \[
k[x,y]/(x,y)^n\longrightarrow H^0(X_n,\mathcal O_{X_n})
\] is an isomorphism. For \(n=1\) both sides are \(k\). At the induction
step, its map on the leftmost ideal layers is \[
(x,y)^{n-1}/(x,y)^n\xrightarrow{\sim}H^0(E,\mathcal O_E(n-1)),
\] the monomial basis identification from projective-space cohomology.
Both the polynomial quotient sequence and the section sequence are short
exact, the latter because the layer has zero \(H^1\). Induction proves
the isomorphism. Their limits give \(k[[x,y]]\). The unit map from the
base local ring to \((b_*\mathcal O_X)_0\) becomes an isomorphism after
completion. Exactness and faithful flatness of completion for finite
modules make it an isomorphism before completion; off the origin the
same holds directly. Thus \(b_*\mathcal O_X=\mathcal O_{\mathbf A^2}\),
as well as higher vanishing.

\subsection{7. Exercises with solutions}\label{exercises-with-solutions}

\textbf{Exercise 7.1 (easy: completed constants).} State and prove the
completed-stalk formula for \(f_*\mathcal O_X\) at a point of a locally
Noetherian base.

\textbf{Solution.} For proper \(f\), Corollary 4.2 in degree zero gives
\(\widehat{(f_*\mathcal O_X)_s}=\varprojlim_n H^0(X_n,\mathcal O_{X_n})\),
with \(X_n\) defined over \(\mathcal O_{S,s}/\mathfrak m_s^n\).
Localization is flat, so identifies the stalk with degree-zero
cohomology over the local ring; global formal functions for that ring
and ideal gives the formula. All maps preserve multiplication, since
they are restriction maps of structure sheaves. Hence it is an
isomorphism of complete algebras, not only of modules.

\textbf{Exercise 7.2 (medium: the exceptional layers).} Compute the
ideal layers of the exceptional curve in the plane blow-up and deduce
\(R^1b_*\mathcal O_X=0\).

\textbf{Solution.} The local ideal frames are \(x\) and \(y\), related
by \(y=vx\); their restriction has the transition of
\(\mathcal O_{\mathbf P^1}(1)\). Since the ideal is invertible, its
\(j\)-th layer is its restricted \(j\)-th tensor power,
\(\mathcal O_{\mathbf P^1}(j)\). Their \(H^1\) groups vanish for
\(j\geq0\), so induction on the exact thickening sequences gives zero
\(H^1\) on every \(X_n\). Formal functions gives zero completed stalk at
the origin, and finite-module Nakayama gives zero stalk. Off the origin
the map is an isomorphism, proving the global assertion.

\textbf{Exercise 7.3 (medium: finite fibers).} Derive finiteness of a
proper morphism with finite fibers using only formal functions, affine
cohomology and Serre's criterion.

\textbf{Solution.} Every fiber and thickening has dimension zero, so its
positive cohomology is zero. Formal functions and Nakayama annihilate
\(R^1f_*J\) for every coherent ideal. Over an affine Noetherian base,
Leray and affine vanishing give \(H^1(X,J)=0\). The finite-ideal affine
criterion makes \(X\) affine. Proper finiteness makes its coordinate
algebra finite over the base algebra, which is exactly finiteness of the
morphism. This also covers empty fibers.

\textbf{Exercise 7.4 (medium: a uniform dimension bound).} If
\(\dim X_s\leq d\) for every \(s\), prove \(R^qf_*F=0\) for \(q>d\). Can
\(\dim X\) replace the fiber calculation in the proof without further
assumptions?

\textbf{Solution.} Theorem 5.1 kills every stalk in the stated degrees,
hence kills the sheaf. The proof works on the Noetherian topology of
each thickened fiber and does not require a dimension bound on the total
space. A locally Noetherian base can have unbounded or infinite Krull
dimension, so a total-space bound might be unavailable and would in any
event miss the sharper relative assertion.

\textbf{Exercise 7.5 (hard: the inverse-limit step).} Prove the
stabilization and limit claims from (1), without assuming the middle
transitions are surjective.

\textbf{Solution.} The right transition \(K_m^{q+1}\to K_n^{q+1}\) is
zero for \(m\geq n+c\), so the middle image lies in \(M^q/J_n^q\). The
left transition is onto, so the middle image contains this entire
submodule. Equality proves Mittag--Leffler stabilization. A compatible
family in the middle maps to a compatible family in the pro-zero right
system. Each of its right coordinates is the image of a sufficiently
distant zero transition, so is zero. Each middle coordinate therefore
has a unique preimage in the injective left term, and uniqueness forces
compatibility of those preimages. Thus the two limits agree. Finally
\(I^nM^q\subset J_n^q\) and \(J_{n+c_J}^q\subset I^nM^q\) identify the
left limit continuously with the adic completion. Every assertion is
justified without finite-level surjectivity of the middle system.

\textbf{Exercise 7.6 (challenging: transient infinitesimal sections).}
For the extension with class \(t\) on \(\mathbf P^1_{k[t]}\), compute
\(H^0\) on every thickening \(B_n=k[t]/(t^n)\), its transition maps, and
its inverse limit.

\textbf{Solution.} Its Grothendieck complex is \([A\xrightarrow{t}A]\)
in degrees zero and one. Over \(B_n\), the kernel of multiplication by
\(t\) is \(k\,t^{n-1}\), including \(k\) when \(n=1\). Reduction
\(B_{n+1}\to B_n\) sends its kernel generator \(t^n\) to zero. Thus all
adjacent transitions on degree-zero cohomology are zero, and its inverse
limit is zero. This agrees with the completion of \(H^0(X,E)=0\),
although every finite thickening has a nonzero section. In degree one,
the cokernel is \(k\) at every level and the transitions are identities;
its limit agrees with the completion of \(A/(t)\). This exhibits both
the transient error and the persistent cohomology in one family.

\subsection{8. Ampleness near a fiber}\label{ampleness-near-a-fiber}

\textbf{Theorem 8.1.} If \(f:X\to S\) is proper with \(S\) locally
Noetherian and \(L|_{X_s}\) is ample, then \(L\) is relatively ample
over an open neighborhood of \(s\).

\textbf{Proof.} Work first over \(A=\mathcal O_{S,s}\), with maximal
ideal \(\mathfrak m\). Put \(X_n=V(\mathfrak m^n\mathcal O_X)\). The
layers \(D_n=\mathfrak m^n\mathcal O_X/\mathfrak m^{n+1}\mathcal O_X\)
form a finite graded module over
\(\mathcal O_{X_1}\otimes_{A/\mathfrak m}\operatorname{gr}_{\mathfrak m}A\):
multiplication gives a surjection from that ring onto their direct sum,
generated in degree zero. The associated graded ring is Noetherian and
finite type. Apply the Serre vanishing theorem to this finite module on
the base change of the proper fiber \(X_1\), with the pulled-back ample
\(L_1\). Affine pushforward and the finite affine-cover complex split
its cohomology by graded degree. Thus one \(d_0\) gives
\(H^1(X_1,D_n\otimes L_1^d)=0\) for every \(n\) and every \(d\ge d_0\).
The exact thickening sequences consequently lift each section of
\(L_1^d\) successively to all \(X_n\). Formal functions identifies the
compatible lift with an element of the completion of the finite module
\(H^0(X,L^d)\). Its reduction modulo \(\mathfrak m\) comes from that
module itself, so the original section lifts to an actual section on
\(X\).

Choose \(d\ge d_0\) for which finitely many sections embed the fiber in
projective space, using the ample embedding lemma in the Serre lesson.
Lift them as above. The local-base cohomology localization proved in the
affine lesson represents the finite collection over an affine
neighborhood of \(s\). Their evaluation cokernel is coherent; its proper
closed support has closed image missing \(s\). Shrink away from that
image. They now generate \(L^d\) and define a morphism
\(h:X\to\mathbf P^r_S\), proper by the closed-graph factorization, whose
fiber at \(s\) is a closed immersion. The quasi-finite open of \(h\)
contains all of \(X_s\). Its closed complement has closed image in
\(S\), since \(X/S\) is proper. Shrinking once more makes \(h\)
quasi-finite everywhere. Theorem 5.2 makes it finite, checked on affine
opens of its target. The inverse images of standard affine projective
charts are now affine; these are exactly the section opens for the
pulled-back coordinate sections of \(L^d=h^*\mathcal O(1)\). They yield
an affine basis after allowing all powers of the line bundle. Indeed, on
one such affine inverse image \(X_{s_i}\), take any principal open
\(D(a)\). Lemma 1.1 of the Serre lesson extends \(a\) to a section \(t\)
of a power of \(L^d\), with \(a=t/s_i^b\) on that chart. Multiplying
\(t\) once more by \(s_i\) forces its nonvanishing open to be contained
in that chart, and that open is exactly \(D(a)\). These principal opens
form a basis. Hence \(L^d\) is relatively ample. The positive-power
criterion proved in the Serre lesson makes \(L\) relatively ample.
\(\square\)

\subsection{References}\label{references}

\begin{itemize}
\tightlist
\item
  \textbf{{[}Stacks{]}} The Stacks project authors, \emph{The Stacks
  project}, in its AI Integrated Stacks Project edition: graded powers
  \href{https://kokunoyumeto.github.io/stacks-zh-hans-cn/en/coherent.html\#coherent-lemma-cohomology-powers-ideal-times-F}{Tag
  02O8}; filtration consequences
  \href{https://kokunoyumeto.github.io/stacks-zh-hans-cn/en/coherent.html\#coherent-lemma-cohomology-powers-ideal-application}{Tag
  02OA}; Mittag--Leffler
  \href{https://kokunoyumeto.github.io/stacks-zh-hans-cn/en/coherent.html\#coherent-lemma-ML-cohomology-powers-ideal}{Tag
  02OB}; formal functions
  \href{https://kokunoyumeto.github.io/stacks-zh-hans-cn/en/coherent.html\#coherent-theorem-formal-functions}{Tag
  02OC}; stalk form
  \href{https://kokunoyumeto.github.io/stacks-zh-hans-cn/en/coherent.html\#coherent-lemma-formal-functions-stalk}{Tag
  02OD}.
\item
  Dimension bound
  \href{https://kokunoyumeto.github.io/stacks-zh-hans-cn/en/coherent.html\#coherent-lemma-higher-direct-images-zero-above-dimension-fibre}{Tag
  02V7}, using the open topological theorem
  \href{https://kokunoyumeto.github.io/stacks-zh-hans-cn/en/cohomology.html\#cohomology-proposition-vanishing-Noetherian}{Tag
  02UZ}; finiteness
  \href{https://kokunoyumeto.github.io/stacks-zh-hans-cn/en/coherent.html\#coherent-lemma-characterize-finite}{Tag
  02OG}. The additional ampleness proof is
  \href{https://kokunoyumeto.github.io/stacks-zh-hans-cn/en/coherent.html\#coherent-lemma-ample-on-fibre}{Tag
  0D2M} and
  \href{https://kokunoyumeto.github.io/stacks-zh-hans-cn/en/coherent.html\#coherent-lemma-ample-in-neighbourhood}{Tag
  0D2N}.
\item
  These open reference treatments retain GNU FDL 1.2. The independently
  authored sections, examples and solutions retain CC0; the adapted
  Appendix Z retains its stated GFDL-1.2-or-later terms. Sections 2--5
  prove the assigned formal-functions, vanishing and finiteness results;
  the completion algebra belongs to the named prerequisite, and Section
  8 proves the additional ampleness consequence.
\end{itemize}

\subsection{Appendix Z. Conductor coefficients and strong
transcendence}\label{appendix-z.-conductor-coefficients-and-strong-transcendence}

This appendix supplies the conductor and strong-transcendence arguments
used in the algebraic Zariski Main proof. The proofs are written here,
at the point where this lesson uses that argument. All rings are
commutative with identity and all ring maps preserve identity.
``Finite'' means finite as a module. The notation \(R[X]\) denotes a
polynomial ring in a formal variable; \(R[x]\subset S\) denotes the
image algebra generated by an element \(x\in S\). No Noetherian
hypothesis is imposed. Reducedness and the domain hypotheses enter only
where stated.

\emph{Adapted from The Stacks Project Authors, Johan de Jong and
contributors, by GPT-6.1 Sol (OpenAI) in Codex at Ultra effort, October
2026. Copyright (C) 2005--2025 Johan de Jong. This adapted appendix is
licensed under the GNU Free Documentation License, version 1.2 or any
later version, with no Invariant Sections, no Front-Cover Texts and no
Back-Cover Texts. A copy of the license is supplied in
\url{COPYING-GFDL-1.2.txt}. The lesson's general CC0 notice does not
replace this appendix's GFDL notice.}

The earlier algebra proofs we use are
\href{../AG-CA/AG-CA-05.html}{Integral extensions: lying over, going up
and going down, Theorems 1.1--1.2 and Proposition 1.3}, for the
determinant test, closure and transitivity of integrality, and finite
base change; its Theorem 3.2 for lying over and Theorem 4.3 for going
down over an arbitrary normal domain;
\href{../AG-CA/AG-CA-02.html}{Localization, local properties and
support, Section 1 and Theorems 3.1--3.2}, for the fraction zero test,
prime correspondence and residue-field fibres;
\href{../AG-CA/AG-CA-01.html}{Spectra of rings, Theorem 1.2, Theorem 5.2
and Corollary 5.3}, for radicals as intersections of primes, clopen
subsets and products; and \href{../AG-CA/AG-CA-03.html}{Noetherian and
Artinian rings, Lemma 4.1 and Theorem 4.2}, for Artinian spectra. The
pointwise quasi-finite test, including finiteness of the residue
extension and of the local fibre algebra, is proved in
\href{../AG-MO/AG-MO-06.html}{Quasi-finite morphisms and Chevalley's
theorem, Theorem 1.1}. Polynomial normality over an arbitrary normal
domain is
\href{the-theorem-on-formal-functions.html\#0-a-normality-input-in-the-algebraic-zariski-main-proof}{Lemma
0.1 of this lesson}.

\subsubsection{Z.1. Integral closure and
localization}\label{z.1.-integral-closure-and-localization}

\textbf{Lemma Z.1.} Let \(A\to B\) be a ring map, let \(T\subset A\) be
multiplicative, and let \(B^{\mathrm{int}}\subset B\) be the subalgebra
of elements integral over \(A\). The integral closure of \(T^{-1}A\) in
\(T^{-1}B\) is \(T^{-1}B^{\mathrm{int}}\), viewed as a subring of
\(T^{-1}B\).

\textbf{Proof.} Localization of the inclusion
\(B^{\mathrm{int}}\subset B\) is still injective: if a fraction from the
first ring becomes zero in the second, a member of \(T\) kills its
numerator already in \(B^{\mathrm{int}}\). A monic equation for
\(b\in B^{\mathrm{int}}\), divided by the appropriate powers of
\(s\in T\), gives a monic equation for \(b/s\) over \(T^{-1}A\). This
proves one inclusion.

For the converse write an integral element as \(z=b/s\), with
\(b\in B\), \(s\in T\), and choose an equation \[
z^d+\sum_{i=1}^d \frac{a_i}{t_i}z^{d-i}=0,
\qquad a_i\in A,\quad t_i\in T.
\] Put \(t=\prod_i t_i\) and \(y=tb\). Multiplication by \((st)^d\)
gives, in \(T^{-1}B\), \[
g(y)=y^d+\sum_{i=1}^d c_i y^{d-i}=0,
\qquad c_i=a_i s^i t^{\,i-1}\prod_{j\ne i}t_j\in A.
\] Here and below coefficients act through the given ring map. The
equation in the localization means that some \(v\in T\) satisfies
\(v g(y)=0\) in \(B\). Consequently \[
(vy)^d+\sum_{i=1}^d v^i c_i (vy)^{d-i}
=v^d g(y)=0
\] in \(B\). This is a monic equation over \(A\), so
\(vy\in B^{\mathrm{int}}\), and \(z=vy/(vts)\) belongs to
\(T^{-1}B^{\mathrm{int}}\). No cancellation or nonzero-divisor
hypothesis was used. \(\square\)

\subsubsection{Z.2. Removing a monic
denominator}\label{z.2.-removing-a-monic-denominator}

\textbf{Lemma Z.2.} Let \(\varphi:R[X]\to S\) be a ring map and let
\(t\in S\) be integral over \(R[X]\). Suppose that for some monic
\(p\in R[X]\) the element \(t\varphi(p)\) belongs to
\(\operatorname{im}\varphi\). There is \(q\in R[X]\) such that
\(t-\varphi(q)\) is integral over \(R\).

\textbf{Proof.} If \(S\) is the zero ring, take \(q=0\), so assume
otherwise. Choose \(r\in R[X]\) with \(t\varphi(p)=\varphi(r)\).
Division by a monic polynomial works over every ring: if the current
dividend has degree at least \(\deg p\), subtract its leading
coefficient times the appropriate power of \(X\) times \(p\); its degree
decreases, and finitely many repetitions give \(r=qp+r'\) with
\(\deg r'<\deg p\). If \(p=1\), take \(q=r\); the difference is zero.
Otherwise put \(\tau=t-\varphi(q)\). Integral elements form a subalgebra
by the earlier Theorem 1.2, so \(\tau\) is still integral over \(R[X]\),
and \[
\tau\varphi(p)=\varphi(r').
\] In \(S_\tau\) this becomes \(\varphi(p)-\tau^{-1}\varphi(r')=0\).
Since \(p\) is monic and \(r'\) has smaller degree, this is a monic
equation for \(x=\varphi(X)\) over the subring
\(D=\varphi(R)[1/\tau]\subset S_\tau\). Thus \(D[x]\) is integral over
\(D\). The original integral equation for \(\tau\), now regarded in
\(S_\tau\), also shows that \(\tau\) is integral over \(D[x]\).
Transitivity, proved in the earlier Theorem 1.2, makes \(\tau\) integral
over \(D\). We can therefore write \[
\tau^d+\sum_{i<d}\left(\sum_{j=0}^{n_i}
\frac{\varphi(b_{ij})}{\tau^j}\right)\tau^i=0
\quad\text{in }S_\tau,
\qquad b_{ij}\in R.
\] Choose an integer \(N\) large enough to clear all displayed negative
powers and to kill the resulting localization error. More explicitly,
first clear those powers; the resulting element of \(S\) becomes zero in
\(S_\tau\), so some further power of \(\tau\) kills it. Combining the
two powers gives \[
\tau^{d+N}+\sum_{i<d}\sum_{j=0}^{n_i}
\varphi(b_{ij})\tau^{i+N-j}=0
\quad\text{in }S.
\] Every lower exponent is nonnegative and strictly less than \(d+N\).
This is a monic equation for \(\tau\) over \(R\), even when \(\tau\) is
a zero divisor. \(\square\)

\subsubsection{Z.3. Clearing a leading
coefficient}\label{z.3.-clearing-a-leading-coefficient}

\textbf{Lemma Z.3.} Let \(\varphi:R[X]\to S\) be a ring map, let
\(t\in S\) be integral over \(R[X]\), and let \[
p=\sum_{i=0}^k a_iX^i\in R[X]
\] satisfy \(t\varphi(p)\in\operatorname{im}\varphi\). Then there are an
integer \(n\geq0\) and \(q\in R[X]\) such that
\(\varphi(a_k)^n t-\varphi(q)\) is integral over \(R\).

\textbf{Proof.} Write \(a=a_k\). If \(a\) is nilpotent in \(R\), choose
a power \(a^n=0\) and take \(q=0\). Otherwise localize \(R\) at \(a\),
and \(S\) at \(\varphi(a)\). The polynomial \(p/a\in R_a[X]\) is monic,
so Lemma Z.2 gives \(q'\in R_a[X]\) for which \(t-\varphi_a(q')\) is
integral over \(R_a\). Choose \(e\geq0\) and \(q_0\in R[X]\) such that
\(q_0/1=a^e q'\). Multiplication by \(a^e\) preserves integrality, since
integral elements form an \(R_a\)-subalgebra. Hence the image in
\(S_{\varphi(a)}\) of \[
\delta=\varphi(a)^e t-\varphi(q_0)
\] is integral over \(R_a\).

By Lemma Z.1, this image can be written as \(c/\varphi(a)^j\), where
\(c\in S\) is integral over \(R\). Equality of these two fractions means
that for some \(\ell\geq0\), \[
\varphi(a)^\ell\bigl(\varphi(a)^j\delta-c\bigr)=0
\quad\text{in }S.
\] Thus \(\varphi(a)^{j+\ell}\delta=\varphi(a)^\ell c\) is integral over
\(R\). Taking \(n=e+j+\ell\) and \(q=a^{j+\ell}q_0\) proves the
assertion. This last step explicitly clears the possible localization
error. \(\square\)

\subsubsection{Z.4. The conductor setup}\label{z.4.-the-conductor-setup}

\textbf{Situation Z.4.} Let \(\varphi:R[X]\to S\) be finite, put
\(A=\varphi(R)\), \(x=\varphi(X)\), and
\(C=A[x]=\operatorname{im}\varphi\). Assume that \(A\) is integrally
closed in \(S\): every element of \(S\) integral over \(A\) belongs to
\(A\). This condition concerns this inclusion; it does not require \(A\)
to be a domain. Define the conductor \[
J=\{g\in S:gS\subset C\}.
\] It is an ideal of \(S\): sums preserve the defining condition, and
\(hgS\subset gS\) for \(h\in S\). Also \(J\subset C\), by evaluating on
\(1\).

Choose \(C\)-module generators \(t_1,\ldots,t_r\) of \(S\), adding \(1\)
if needed. For \(g\in S\), we have \(g\in J\) exactly when \(gt_i\in C\)
for all \(i\). Indeed, this condition makes every \(C\)-linear
combination of the \(gt_i\) belong to \(C\). These observations also
cover the zero ring.

\subsubsection{Z.5. The leading coefficient enters the
conductor}\label{z.5.-the-leading-coefficient-enters-the-conductor}

\textbf{Lemma Z.5.} In Situation Z.4, let \(u\in S\) and
\(p=\sum_{i=0}^k a_iX^i\in R[X]\). If \(u\varphi(p)\in J\), there is
\(m\geq0\) such that \(u\varphi(a_k)^m\in J\).

\textbf{Proof.} For each module generator \(t_i\), the element
\(ut_i\in S\) is integral over \(R[X]\), because \(S\) is finite over
\(R[X]\) and the earlier determinant test proves finite algebras
integral. Moreover \[
(ut_i)\varphi(p)=(u\varphi(p))t_i\in C.
\] Lemma Z.3 gives \(n_i\geq0\) and \(q_i\in R[X]\) such that \[
\varphi(a_k)^{n_i}ut_i-\varphi(q_i)
\] is integral over \(R\), equivalently over its image \(A\). The
integral-closedness assumption puts this difference in \(A\subset C\);
hence \(\varphi(a_k)^{n_i}ut_i\in C\). With \(m=\max_i n_i\),
multiplication by additional powers of \(\varphi(a_k)\in C\) gives
\(u\varphi(a_k)^m t_i\in C\) for every \(i\). The generator
characterization in Situation Z.4 now gives \(u\varphi(a_k)^m\in J\).
\(\square\)

\subsubsection{Z.6. Every coefficient enters the
radical}\label{z.6.-every-coefficient-enters-the-radical}

\textbf{Lemma Z.6.} In Situation Z.4, if \(u\varphi(p)\in\sqrt J\),
where \(u\in S\) and \(p=\sum_{i=0}^k a_iX^i\), then
\(u\varphi(a_i)\in\sqrt J\) for every \(i\).

\textbf{Proof.} Choose \(N\geq1\) with \(u^N\varphi(p^N)\in J\). The
coefficient of \(X^{kN}\) in \(p^N\) is \(a_k^N\), so Lemma Z.5 gives \[
u^N\varphi(a_k)^{Nm}\in J
\] for some \(m\geq0\). Increase \(m\) to at least \(1\), which is
allowed because \(J\) is an ideal. Then \[
(u\varphi(a_k))^{Nm}
=u^{N(m-1)}\bigl(u^N\varphi(a_k)^{Nm}\bigr)\in J.
\] Thus \(u\varphi(a_k)\in\sqrt J\). Subtracting the term
\(u\varphi(a_k)x^k\) from \(u\varphi(p)\) leaves
\(u\varphi(p-a_kX^k)\in\sqrt J\). Induction on the chosen top index
\(k\) proves the assertion for all coefficients; for \(k=0\) the first
step already proves it. \(\square\)

\subsubsection{Z.7. Strong transcendence and the conductor
quotient}\label{z.7.-strong-transcendence-and-the-conductor-quotient}

\textbf{Definition Z.7.} For an inclusion \(R\subset S\), an element
\(x\in S\) is \textbf{strongly transcendental over \(R\)} if every
equation \[
u\sum_{i=0}^k a_ix^i=0,
\qquad u\in S,\quad a_i\in R,
\] implies \(ua_i=0\) for every \(i\). If \(R\subset S\) are domains,
this is equivalent to ordinary transcendence: the case \(u=1\) gives
injectivity of \(R[X]\to S\), and conversely a nonzero \(u\) can be
cancelled in a domain.

In Situation Z.4 put \[
S_0=S/\sqrt J,\qquad A_0=A/(A\cap\sqrt J)\subset S_0,
\] and let \(x_0\) be the image of \(x\). Then \(A_0\) and \(S_0\) are
reduced, \(S_0\) is finite over \(A_0[x_0]\), and \(x_0\) is strongly
transcendental over \(A_0\).

\textbf{Proof of the assertion.} The displayed map from \(A_0\) is
injective by its definition. The ring \(S_0\) is reduced because the
quotient ideal is radical; its subring \(A_0\) is therefore reduced. The
images of the finite \(C\)-module generators of \(S\) generate \(S_0\)
over the image \(A_0[x_0]\).

Given a relation \(u_0\sum_i\bar a_i x_0^i=0\), lift \(u_0\) to
\(u\in S\), the \(\bar a_i\) to \(a_i\in A\), and these coefficients to
\(r_i\in R\) with \(\varphi(r_i)=a_i\). The relation says
\(u\varphi(\sum_i r_iX^i)\in\sqrt J\). Lemma Z.6 gives
\(ua_i\in\sqrt J\) for every \(i\), which is precisely \(u_0\bar a_i=0\)
in \(S_0\). \(\square\)

\subsubsection{Z.8. Quasi-finiteness in a four-ring
diagram}\label{z.8.-quasi-finiteness-in-a-four-ring-diagram}

\textbf{Lemma Z.8.} Consider a commutative square of rings \[
\begin{array}{ccc}
R&\longrightarrow&S\\
\downarrow&&\downarrow\\
R'&\longrightarrow&S'.
\end{array}
\] Let \(\mathfrak q'\subset S'\) contract to \(\mathfrak q\subset S\)
and \(\mathfrak p'\subset R'\); write
\(\mathfrak p=\mathfrak q\cap R=\mathfrak p'\cap R\). Suppose \(R\to S\)
is of finite type and the canonical map \(S\otimes_R R'\to S'\) is
surjective. If \(R\to S\) is quasi-finite at \(\mathfrak q\), then
\(R'\to S'\) is quasi-finite at \(\mathfrak q'\).

\textbf{Proof.} Put \(k=\kappa(\mathfrak p)\),
\(k'=\kappa(\mathfrak p')\), \(F=S\otimes_R k\), and
\(F'=S'\otimes_{R'}k'\). The contractions give field inclusions
\(k\subset k'\) and \(\kappa(\mathfrak q)\subset\kappa(\mathfrak q')\).
The fibre prime corresponding to \(\mathfrak q\) is isolated by
quasi-finiteness, and is closed by the earlier pointwise test. Its
singleton is consequently clopen in \(\operatorname{Spec}F\).

The clopen-idempotent theorem and product decomposition proved in the
earlier spectra lesson give \[
F=F_1\times F_2,
\] with the singleton as \(\operatorname{Spec}F_1\). The ring \(F_1\)
has only one prime, which is its unique maximal ideal. Elements outside
this ideal are units, so localization at it does not change \(F_1\).
Thus \(F_1\) is the local fibre ring at \(\mathfrak q\), and the earlier
pointwise test makes \(F_1\) finite-dimensional over \(k\).

Tensoring the assumed surjection with \(k'\) gives a surjection \[
F\otimes_k k'\longrightarrow F'.
\] The complementary idempotents of \(F\) give \(F'=F'_1\times F'_2\),
with surjections \(F_i\otimes_k k'\to F'_i\). In particular \(F'_1\) is
finite-dimensional over \(k'\). The composite
\(F\to F'\to\kappa(\mathfrak q')\) factors through
\(\kappa(\mathfrak q)\): it is induced by the ring map and the
contraction \(\mathfrak q'\cap S=\mathfrak q\). Therefore it kills the
\(F_2\) idempotent, just as \(F\to\kappa(\mathfrak q)\) does. The point
\(\mathfrak q'\) consequently belongs to \(\operatorname{Spec}F'_1\).

A finite-dimensional algebra over a field is Artinian, since a
descending chain of ideals is a descending chain of vector subspaces.
The earlier Artinian theorem makes its spectrum a finite set of closed
points, hence discrete: the complement of each point is a finite union
of closed points. Thus \(\mathfrak q'\) is isolated in
\(\operatorname{Spec}F'_1\), and in \(\operatorname{Spec}F'\), because
this factor is open and closed. Finally \(R'\to S'\) is of finite type:
\(S'\) is a quotient of the finite type \(R'\)-algebra
\(S\otimes_R R'\). This verifies both parts of the quasi-finite
condition. \(\square\)

\subsubsection{Z.9. Passage to a minimal
prime}\label{z.9.-passage-to-a-minimal-prime}

\textbf{Lemma Z.9.} Let \(R\subset S\) be an inclusion of reduced rings,
let \(x\in S\) be strongly transcendental over \(R\), and let
\(\mathfrak q\) be a minimal prime of \(S\). Put
\(\mathfrak p=R\cap\mathfrak q\). The image of \(x\) in
\(S/\mathfrak q\) is strongly transcendental over \(R/\mathfrak p\).

\textbf{Proof.} We first justify the local fact used here without a
Noetherian assumption. The localization \(S_{\mathfrak q}\) is reduced.
Indeed, if \((y/s)^n=0\), then some \(v\notin\mathfrak q\) satisfies
\(vy^n=0\); hence \((vy)^n=0\) and reducedness of \(S\) gives \(vy=0\),
so \(y/s=0\). Prime correspondence identifies the primes of
\(S_{\mathfrak q}\) with primes of \(S\) contained in \(\mathfrak q\).
Minimality leaves only \(\mathfrak qS_{\mathfrak q}\). The
intersection-of-primes description of the nilradical, together with
reducedness, therefore gives \(\mathfrak qS_{\mathfrak q}=0\). This
local ring has zero maximal ideal, so is a field.

Now suppose \(u\sum_i a_ix^i\in\mathfrak q\), with \(u\in S\) and
\(a_i\in R\). This element becomes zero in \(S_{\mathfrak q}\), so there
is \(v\notin\mathfrak q\) such that \(vu\sum_i a_ix^i=0\) in \(S\).
Strong transcendence gives \(vua_i=0\) for every \(i\). Since
\(\mathfrak q\) is prime and \(v\notin\mathfrak q\), we get
\(ua_i\in\mathfrak q\) for every \(i\). These are exactly the required
coefficient equalities after quotienting. The induced map
\(R/\mathfrak p\to S/\mathfrak q\) is an inclusion by the contraction
defining \(\mathfrak p\). \(\square\)

\subsubsection{Z.10. Transcendence excludes quasi-finite points for
domains}\label{z.10.-transcendence-excludes-quasi-finite-points-for-domains}

\textbf{Lemma Z.10.} Suppose \(R\subset S\) are domains, \(x\in S\) is
transcendental over \(R\), and \(S\) is finite over \(R[x]\). Then
\(R\to S\) is not quasi-finite at any prime of \(S\).

\textbf{Proof.} The map \(R\to S\) is of finite type: the element \(x\),
together with a finite list of \(R[x]\)-module generators of \(S\),
generates \(S\) as an \(R\)-algebra.

First assume \(R\) is normal, meaning integrally closed in its fraction
field. Transcendence identifies \(R[x]\) with \(R[X]\); Lemma 0.1 above
proves that it is normal. Fix \(\mathfrak q\subset S\), put
\(\mathfrak p=R\cap\mathfrak q\) and
\(\mathfrak r=R[x]\cap\mathfrak q\), and suppose for contradiction that
\(R\to S\) is quasi-finite at \(\mathfrak q\). The earlier pointwise
test gives a finite field extension
\(\kappa(\mathfrak p)\subset\kappa(\mathfrak q)\). The intermediate
field \[
\kappa(\mathfrak p)\subset\kappa(\mathfrak r)
\subset\kappa(\mathfrak q)
\] is therefore finite over \(\kappa(\mathfrak p)\). The containment
\(\mathfrak pR[x]\subset\mathfrak r\) must be strict: equality would
give \(\kappa(\mathfrak r)=\kappa(\mathfrak p)(X)\), a transcendental
extension.

The inclusion \(R[x]\subset S\) is integral, by its finiteness. The
earlier going-down theorem for an integral inclusion of domains with
normal base gives a prime \(\mathfrak q_0\subset\mathfrak q\) of \(S\)
contracting to \(\mathfrak pR[x]\). It is distinct from \(\mathfrak q\),
because their contractions to \(R[x]\) differ. Both primes contract to
\(\mathfrak p\) in \(R\), so they are distinct points of the same fibre,
with \(\mathfrak q_0\) specializing to \(\mathfrak q\). Every open
containing a point contains its generizations, as is seen on the basic
opens \(D(f)\). Thus \(\mathfrak q\) cannot be isolated in that fibre.
This contradicts quasi-finiteness.

For arbitrary \(R\), put \(K=\operatorname{Frac}R\),
\(L=\operatorname{Frac}S\), let \(R'\) be the integral closure of \(R\)
in \(K\), and let \(S'\subset L\) be the subring generated by \(S\) and
\(R'\). These are domains. The ring \(R'\) is normal: an element of
\(K\) integral over \(R'\) is integral over \(R\) by transitivity, and
hence belongs to \(R'\). Its fraction field is \(K\), because
\(R\subset R'\subset K\).

The element \(x\) remains transcendental over \(R'\). Any relation over
\(R'\subset K\), after multiplying by a common nonzero denominator from
\(R\), would give a relation over \(R\). By construction
\(S\otimes_R R'\to S'\) is surjective. If \(t_1,\ldots,t_r\) generate
\(S\) over \(R[x]\), their images generate \(S'\) over \(R'[x]\): they
generate the tensor product over \(R[x]\otimes_R R'=R'[X]\), and
therefore its quotient \(S'\). Hence \(S'\) is finite over \(R'[x]\).

Every element of \(R'\) is integral over \(S\), since its monic equation
has coefficients in \(R\subset S\). Closure of integral elements under
algebra generation makes \(S\subset S'\) integral. Lying over, in the
earlier Theorem 3.2, gives a prime \(\mathfrak q'\subset S'\) above any
prescribed prime \(\mathfrak q\subset S\). If \(R\to S\) were
quasi-finite at \(\mathfrak q\), Lemma Z.8 applied to \(R,S,R',S'\)
would make \(R'\to S'\) quasi-finite at \(\mathfrak q'\), contradicting
the normal case. This uses neither finiteness of \(R'\) over \(R\) nor
Noetherianity. \(\square\)

\subsubsection{Z.11. Strong transcendence excludes quasi-finite
points}\label{z.11.-strong-transcendence-excludes-quasi-finite-points}

\textbf{Lemma Z.11.} Suppose \(R\subset S\) are reduced rings,
\(x\in S\) is strongly transcendental over \(R\), and \(S\) is finite
over \(R[x]\). Then \(R\to S\) is not quasi-finite at any prime of
\(S\).

\textbf{Proof.} Again \(R\to S\) is of finite type. Fix a prime
\(\mathfrak q\subset S\) and choose a minimal prime
\(\mathfrak q_0\subset\mathfrak q\). Such a prime exists without
Noetherianity: order the primes contained in \(\mathfrak q\) by reverse
inclusion. A nonempty chain has an upper bound given by its
intersection, which is a proper prime ideal contained in
\(\mathfrak q\). To check primality, if neither \(a\) nor \(b\) belongs
to that intersection, choose members of the chain omitting each; the
smaller of these two prime ideals omits both and therefore omits \(ab\).
The empty chain is bounded by \(\mathfrak q\). Zorn's lemma now gives a
minimal prime.

Put \(\mathfrak p_0=R\cap\mathfrak q_0\), \(R_0=R/\mathfrak p_0\), and
\(S_0=S/\mathfrak q_0\). The induced inclusion \(R_0\subset S_0\) is an
inclusion of domains. By Lemma Z.9, the image \(x_0\) is strongly
transcendental, and hence transcendental, over \(R_0\). The quotient
\(S_0\) remains finite over \(R_0[x_0]\), by the images of the same
module generators. Lemma Z.10 says that \(R_0\to S_0\) is not
quasi-finite at \(\mathfrak q/\mathfrak q_0\).

The square \(R,S,R_0,S_0\) satisfies the hypotheses of Lemma Z.8: its
tensor-product map is surjective, because
\(S\otimes_R R_0=S/\mathfrak p_0S\) surjects onto \(S/\mathfrak q_0\).
If \(R\to S\) were quasi-finite at \(\mathfrak q\), that lemma would
make \(R_0\to S_0\) quasi-finite at \(\mathfrak q/\mathfrak q_0\), a
contradiction. This proof does not require finitely many minimal primes.
\(\square\)

\subsubsection{Z.12. The conductor is avoided at a quasi-finite
point}\label{z.12.-the-conductor-is-avoided-at-a-quasi-finite-point}

\textbf{Corollary Z.12.} In Situation Z.4, if \(A\to S\) is quasi-finite
at a prime \(\mathfrak q\subset S\), then \(J\not\subset\mathfrak q\).

\textbf{Proof.} Suppose \(J\subset\mathfrak q\). Then
\(\sqrt J\subset\mathfrak q\). Use the reduced rings
\(A_0\subset S_0=S/\sqrt J\) and the strongly transcendental element
\(x_0\) constructed in Z.7. They satisfy the hypotheses of Lemma Z.11,
so \(A_0\to S_0\) is not quasi-finite at \(\mathfrak q/\sqrt J\).

On the other hand \(A\to S\) is of finite type, because \(S\) is finite
over \(A[x]\). The square \(A,S,A_0,S_0\) has surjective tensor-product
map: the quotient \(S/(A\cap\sqrt J)S\) surjects onto \(S/\sqrt J\).
Lemma Z.8 therefore transfers quasi-finiteness at \(\mathfrak q\) to
quasi-finiteness at \(\mathfrak q/\sqrt J\), giving the contradiction.
\(\square\)

For the finite-over-one-generator step of
\href{../AG-MO/AG-MO-12.html}{Zariski's Main Theorem, Section 2, Lemma
2.1}, this corollary supplies precisely the conductor element outside
the chosen prime. Lemma Z.6 supplies the preceding coefficient
assertion, and Lemma Z.11 supplies the reduced strong-transcendence
assertion. These are the two supporting arguments required there; the
earlier lesson contains the separate monogenic argument and the main
induction.

\textbf{Source and license record.} The arguments adapted here are the
bodies with Stacks tags 0307, 00PT, 00PV, 00PW (the situation), 00PX,
00PY, 00PN, 00Q0, 00Q1 and 00Q2, together with the definition labelled
``definition-strongly-transcendental'', in the
\href{https://github.com/KokunoYumeto/unofficial-stacks-project-ai-drafts/blob/565b10e987aba5969b21145a0833f42d69f96790/algebra.tex}{frozen
AI Integrated Stacks Project algebra source}. Tags identify the adapted
bodies and their authorship; the mathematical dependencies in this
appendix are the written proofs above and the exact earlier programme
lessons identified in its opening paragraph. The source's GFDL
authorship and license conditions remain attached to this adapted text.
\clearpage\section*{Copyright and licence of Stacks adaptations}
The identified Stacks proof adaptations retain the authorship of the Stacks project authors and AI Integrated Stacks Project contributors. Adaptations and added exposition by GPT-6.1 Sol (OpenAI), Codex, Ultra, 5 October 2026. Permission is granted to copy, distribute and/or modify these identified components under the GNU Free Documentation License, Version 1.2 or any later version published by the Free Software Foundation; with no Invariant Sections, no Front-Cover Texts, and no Back-Cover Texts. Independently authored portions retain their stated CC0 terms. A copy of the licence follows.

\begingroup\footnotesize
\begin{verbatim}
		GNU Free Documentation License
		  Version 1.2, November 2002


 Copyright (C) 2000,2001,2002  Free Software Foundation, Inc.
     51 Franklin St, Fifth Floor, Boston, MA  02110-1301  USA
 Everyone is permitted to copy and distribute verbatim copies
 of this license document, but changing it is not allowed.


0. PREAMBLE

The purpose of this License is to make a manual, textbook, or other
functional and useful document "free" in the sense of freedom: to
assure everyone the effective freedom to copy and redistribute it,
with or without modifying it, either commercially or noncommercially.
Secondarily, this License preserves for the author and publisher a way
to get credit for their work, while not being considered responsible
for modifications made by others.

This License is a kind of "copyleft", which means that derivative
works of the document must themselves be free in the same sense.  It
complements the GNU General Public License, which is a copyleft
license designed for free software.

We have designed this License in order to use it for manuals for free
software, because free software needs free documentation: a free
program should come with manuals providing the same freedoms that the
software does.  But this License is not limited to software manuals;
it can be used for any textual work, regardless of subject matter or
whether it is published as a printed book.  We recommend this License
principally for works whose purpose is instruction or reference.


1. APPLICABILITY AND DEFINITIONS

This License applies to any manual or other work, in any medium, that
contains a notice placed by the copyright holder saying it can be
distributed under the terms of this License.  Such a notice grants a
world-wide, royalty-free license, unlimited in duration, to use that
work under the conditions stated herein.  The "Document", below,
refers to any such manual or work.  Any member of the public is a
licensee, and is addressed as "you".  You accept the license if you
copy, modify or distribute the work in a way requiring permission
under copyright law.

A "Modified Version" of the Document means any work containing the
Document or a portion of it, either copied verbatim, or with
modifications and/or translated into another language.

A "Secondary Section" is a named appendix or a front-matter section of
the Document that deals exclusively with the relationship of the
publishers or authors of the Document to the Document's overall subject
(or to related matters) and contains nothing that could fall directly
within that overall subject.  (Thus, if the Document is in part a
textbook of mathematics, a Secondary Section may not explain any
mathematics.)  The relationship could be a matter of historical
connection with the subject or with related matters, or of legal,
commercial, philosophical, ethical or political position regarding
them.

The "Invariant Sections" are certain Secondary Sections whose titles
are designated, as being those of Invariant Sections, in the notice
that says that the Document is released under this License.  If a
section does not fit the above definition of Secondary then it is not
allowed to be designated as Invariant.  The Document may contain zero
Invariant Sections.  If the Document does not identify any Invariant
Sections then there are none.

The "Cover Texts" are certain short passages of text that are listed,
as Front-Cover Texts or Back-Cover Texts, in the notice that says that
the Document is released under this License.  A Front-Cover Text may
be at most 5 words, and a Back-Cover Text may be at most 25 words.

A "Transparent" copy of the Document means a machine-readable copy,
represented in a format whose specification is available to the
general public, that is suitable for revising the document
straightforwardly with generic text editors or (for images composed of
pixels) generic paint programs or (for drawings) some widely available
drawing editor, and that is suitable for input to text formatters or
for automatic translation to a variety of formats suitable for input
to text formatters.  A copy made in an otherwise Transparent file
format whose markup, or absence of markup, has been arranged to thwart
or discourage subsequent modification by readers is not Transparent.
An image format is not Transparent if used for any substantial amount
of text.  A copy that is not "Transparent" is called "Opaque".

Examples of suitable formats for Transparent copies include plain
ASCII without markup, Texinfo input format, LaTeX input format, SGML
or XML using a publicly available DTD, and standard-conforming simple
HTML, PostScript or PDF designed for human modification.  Examples of
transparent image formats include PNG, XCF and JPG.  Opaque formats
include proprietary formats that can be read and edited only by
proprietary word processors, SGML or XML for which the DTD and/or
processing tools are not generally available, and the
machine-generated HTML, PostScript or PDF produced by some word
processors for output purposes only.

The "Title Page" means, for a printed book, the title page itself,
plus such following pages as are needed to hold, legibly, the material
this License requires to appear in the title page.  For works in
formats which do not have any title page as such, "Title Page" means
the text near the most prominent appearance of the work's title,
preceding the beginning of the body of the text.

A section "Entitled XYZ" means a named subunit of the Document whose
title either is precisely XYZ or contains XYZ in parentheses following
text that translates XYZ in another language.  (Here XYZ stands for a
specific section name mentioned below, such as "Acknowledgements",
"Dedications", "Endorsements", or "History".)  To "Preserve the Title"
of such a section when you modify the Document means that it remains a
section "Entitled XYZ" according to this definition.

The Document may include Warranty Disclaimers next to the notice which
states that this License applies to the Document.  These Warranty
Disclaimers are considered to be included by reference in this
License, but only as regards disclaiming warranties: any other
implication that these Warranty Disclaimers may have is void and has
no effect on the meaning of this License.


2. VERBATIM COPYING

You may copy and distribute the Document in any medium, either
commercially or noncommercially, provided that this License, the
copyright notices, and the license notice saying this License applies
to the Document are reproduced in all copies, and that you add no other
conditions whatsoever to those of this License.  You may not use
technical measures to obstruct or control the reading or further
copying of the copies you make or distribute.  However, you may accept
compensation in exchange for copies.  If you distribute a large enough
number of copies you must also follow the conditions in section 3.

You may also lend copies, under the same conditions stated above, and
you may publicly display copies.


3. COPYING IN QUANTITY

If you publish printed copies (or copies in media that commonly have
printed covers) of the Document, numbering more than 100, and the
Document's license notice requires Cover Texts, you must enclose the
copies in covers that carry, clearly and legibly, all these Cover
Texts: Front-Cover Texts on the front cover, and Back-Cover Texts on
the back cover.  Both covers must also clearly and legibly identify
you as the publisher of these copies.  The front cover must present
the full title with all words of the title equally prominent and
visible.  You may add other material on the covers in addition.
Copying with changes limited to the covers, as long as they preserve
the title of the Document and satisfy these conditions, can be treated
as verbatim copying in other respects.

If the required texts for either cover are too voluminous to fit
legibly, you should put the first ones listed (as many as fit
reasonably) on the actual cover, and continue the rest onto adjacent
pages.

If you publish or distribute Opaque copies of the Document numbering
more than 100, you must either include a machine-readable Transparent
copy along with each Opaque copy, or state in or with each Opaque copy
a computer-network location from which the general network-using
public has access to download using public-standard network protocols
a complete Transparent copy of the Document, free of added material.
If you use the latter option, you must take reasonably prudent steps,
when you begin distribution of Opaque copies in quantity, to ensure
that this Transparent copy will remain thus accessible at the stated
location until at least one year after the last time you distribute an
Opaque copy (directly or through your agents or retailers) of that
edition to the public.

It is requested, but not required, that you contact the authors of the
Document well before redistributing any large number of copies, to give
them a chance to provide you with an updated version of the Document.


4. MODIFICATIONS

You may copy and distribute a Modified Version of the Document under
the conditions of sections 2 and 3 above, provided that you release
the Modified Version under precisely this License, with the Modified
Version filling the role of the Document, thus licensing distribution
and modification of the Modified Version to whoever possesses a copy
of it.  In addition, you must do these things in the Modified Version:

A. Use in the Title Page (and on the covers, if any) a title distinct
   from that of the Document, and from those of previous versions
   (which should, if there were any, be listed in the History section
   of the Document).  You may use the same title as a previous version
   if the original publisher of that version gives permission.
B. List on the Title Page, as authors, one or more persons or entities
   responsible for authorship of the modifications in the Modified
   Version, together with at least five of the principal authors of the
   Document (all of its principal authors, if it has fewer than five),
   unless they release you from this requirement.
C. State on the Title page the name of the publisher of the
   Modified Version, as the publisher.
D. Preserve all the copyright notices of the Document.
E. Add an appropriate copyright notice for your modifications
   adjacent to the other copyright notices.
F. Include, immediately after the copyright notices, a license notice
   giving the public permission to use the Modified Version under the
   terms of this License, in the form shown in the Addendum below.
G. Preserve in that license notice the full lists of Invariant Sections
   and required Cover Texts given in the Document's license notice.
H. Include an unaltered copy of this License.
I. Preserve the section Entitled "History", Preserve its Title, and add
   to it an item stating at least the title, year, new authors, and
   publisher of the Modified Version as given on the Title Page.  If
   there is no section Entitled "History" in the Document, create one
   stating the title, year, authors, and publisher of the Document as
   given on its Title Page, then add an item describing the Modified
   Version as stated in the previous sentence.
J. Preserve the network location, if any, given in the Document for
   public access to a Transparent copy of the Document, and likewise
   the network locations given in the Document for previous versions
   it was based on.  These may be placed in the "History" section.
   You may omit a network location for a work that was published at
   least four years before the Document itself, or if the original
   publisher of the version it refers to gives permission.
K. For any section Entitled "Acknowledgements" or "Dedications",
   Preserve the Title of the section, and preserve in the section all
   the substance and tone of each of the contributor acknowledgements
   and/or dedications given therein.
L. Preserve all the Invariant Sections of the Document,
   unaltered in their text and in their titles.  Section numbers
   or the equivalent are not considered part of the section titles.
M. Delete any section Entitled "Endorsements".  Such a section
   may not be included in the Modified Version.
N. Do not retitle any existing section to be Entitled "Endorsements"
   or to conflict in title with any Invariant Section.
O. Preserve any Warranty Disclaimers.

If the Modified Version includes new front-matter sections or
appendices that qualify as Secondary Sections and contain no material
copied from the Document, you may at your option designate some or all
of these sections as invariant.  To do this, add their titles to the
list of Invariant Sections in the Modified Version's license notice.
These titles must be distinct from any other section titles.

You may add a section Entitled "Endorsements", provided it contains
nothing but endorsements of your Modified Version by various
parties--for example, statements of peer review or that the text has
been approved by an organization as the authoritative definition of a
standard.

You may add a passage of up to five words as a Front-Cover Text, and a
passage of up to 25 words as a Back-Cover Text, to the end of the list
of Cover Texts in the Modified Version.  Only one passage of
Front-Cover Text and one of Back-Cover Text may be added by (or
through arrangements made by) any one entity.  If the Document already
includes a cover text for the same cover, previously added by you or
by arrangement made by the same entity you are acting on behalf of,
you may not add another; but you may replace the old one, on explicit
permission from the previous publisher that added the old one.

The author(s) and publisher(s) of the Document do not by this License
give permission to use their names for publicity for or to assert or
imply endorsement of any Modified Version.


5. COMBINING DOCUMENTS

You may combine the Document with other documents released under this
License, under the terms defined in section 4 above for modified
versions, provided that you include in the combination all of the
Invariant Sections of all of the original documents, unmodified, and
list them all as Invariant Sections of your combined work in its
license notice, and that you preserve all their Warranty Disclaimers.

The combined work need only contain one copy of this License, and
multiple identical Invariant Sections may be replaced with a single
copy.  If there are multiple Invariant Sections with the same name but
different contents, make the title of each such section unique by
adding at the end of it, in parentheses, the name of the original
author or publisher of that section if known, or else a unique number.
Make the same adjustment to the section titles in the list of
Invariant Sections in the license notice of the combined work.

In the combination, you must combine any sections Entitled "History"
in the various original documents, forming one section Entitled
"History"; likewise combine any sections Entitled "Acknowledgements",
and any sections Entitled "Dedications".  You must delete all sections
Entitled "Endorsements".


6. COLLECTIONS OF DOCUMENTS

You may make a collection consisting of the Document and other documents
released under this License, and replace the individual copies of this
License in the various documents with a single copy that is included in
the collection, provided that you follow the rules of this License for
verbatim copying of each of the documents in all other respects.

You may extract a single document from such a collection, and distribute
it individually under this License, provided you insert a copy of this
License into the extracted document, and follow this License in all
other respects regarding verbatim copying of that document.


7. AGGREGATION WITH INDEPENDENT WORKS

A compilation of the Document or its derivatives with other separate
and independent documents or works, in or on a volume of a storage or
distribution medium, is called an "aggregate" if the copyright
resulting from the compilation is not used to limit the legal rights
of the compilation's users beyond what the individual works permit.
When the Document is included in an aggregate, this License does not
apply to the other works in the aggregate which are not themselves
derivative works of the Document.

If the Cover Text requirement of section 3 is applicable to these
copies of the Document, then if the Document is less than one half of
the entire aggregate, the Document's Cover Texts may be placed on
covers that bracket the Document within the aggregate, or the
electronic equivalent of covers if the Document is in electronic form.
Otherwise they must appear on printed covers that bracket the whole
aggregate.


8. TRANSLATION

Translation is considered a kind of modification, so you may
distribute translations of the Document under the terms of section 4.
Replacing Invariant Sections with translations requires special
permission from their copyright holders, but you may include
translations of some or all Invariant Sections in addition to the
original versions of these Invariant Sections.  You may include a
translation of this License, and all the license notices in the
Document, and any Warranty Disclaimers, provided that you also include
the original English version of this License and the original versions
of those notices and disclaimers.  In case of a disagreement between
the translation and the original version of this License or a notice
or disclaimer, the original version will prevail.

If a section in the Document is Entitled "Acknowledgements",
"Dedications", or "History", the requirement (section 4) to Preserve
its Title (section 1) will typically require changing the actual
title.


9. TERMINATION

You may not copy, modify, sublicense, or distribute the Document except
as expressly provided for under this License.  Any other attempt to
copy, modify, sublicense or distribute the Document is void, and will
automatically terminate your rights under this License.  However,
parties who have received copies, or rights, from you under this
License will not have their licenses terminated so long as such
parties remain in full compliance.


10. FUTURE REVISIONS OF THIS LICENSE

The Free Software Foundation may publish new, revised versions
of the GNU Free Documentation License from time to time.  Such new
versions will be similar in spirit to the present version, but may
differ in detail to address new problems or concerns.  See
https://www.gnu.org/copyleft/.

Each version of the License is given a distinguishing version number.
If the Document specifies that a particular numbered version of this
License "or any later version" applies to it, you have the option of
following the terms and conditions either of that specified version or
of any later version that has been published (not as a draft) by the
Free Software Foundation.  If the Document does not specify a version
number of this License, you may choose any version ever published (not
as a draft) by the Free Software Foundation.


ADDENDUM: How to use this License for your documents

To use this License in a document you have written, include a copy of
the License in the document and put the following copyright and
license notices just after the title page:

    Copyright (c)  YEAR  YOUR NAME.
    Permission is granted to copy, distribute and/or modify this document
    under the terms of the GNU Free Documentation License, Version 1.2
    or any later version published by the Free Software Foundation;
    with no Invariant Sections, no Front-Cover Texts, and no Back-Cover Texts.
    A copy of the license is included in the section entitled "GNU
    Free Documentation License".

If you have Invariant Sections, Front-Cover Texts and Back-Cover Texts,
replace the "with...Texts." line with this:

    with the Invariant Sections being LIST THEIR TITLES, with the
    Front-Cover Texts being LIST, and with the Back-Cover Texts being LIST.

If you have Invariant Sections without Cover Texts, or some other
combination of the three, merge those two alternatives to suit the
situation.

If your document contains nontrivial examples of program code, we
recommend releasing these examples in parallel under your choice of
free software license, such as the GNU General Public License,
to permit their use in free software.

\end{verbatim}\endgroup
\end{document}
